\section{Introduction}
\label{sec:introduction}

Large Language Models (LLMs) have moved Natural Language Processing (NLP) from task-specific classifiers towards general-purpose generative systems that can follow instructions, imitate styles and write fluent text~\cite{brown2020language, openai2023gpt4}. This creates a dual-use situation for computational text and media sciences. On the one hand, LLMs are studied as \emph{synthetic participants} that might reproduce the behaviour or opinions of human groups~\cite{aher2023using, santurkar2023whose}. On the other hand, the same fluency lowers the cost of producing misleading or manipulative text~\cite{wardle2017information, buchanan2021truth, vykopal2024disinformation}, which makes automated detection of such text more important~\cite{dasanmartino2020semeval, gaeta2025towards}.

A central question behind both uses is representational fidelity. Instruction tuning and reinforcement learning from human feedback (RLHF) make models more helpful and compliant~\cite{ouyang2022training}, but they do not make them a neutral mirror of the population. Santurkar et al.\ find that the opinions expressed by LMs are substantially misaligned with those of US demographic groups, and that human-feedback-tuned models shift towards more liberal, educated and wealthy groups~\cite{santurkar2023whose}. At the same time, LLMs are increasingly used as tools for detecting propaganda techniques in news text~\cite{gaeta2025towards}.

This term paper combines a structured comparative review of four recent empirical studies with a small experiment of its own. The studies each cover one side of the dual role: Aher et al.~\cite{aher2023using} (simulating human subjects), Santurkar et al.~\cite{santurkar2023whose} (whose opinions LMs reflect), Vykopal et al.~\cite{vykopal2024disinformation} (disinformation generation) and Gaeta et al.~\cite{gaeta2025towards} (LLM-based propaganda detection). The paper is guided by four research questions; RQ1 is additionally tested in the experiment:
\begin{enumerate}
    \item \textbf{RQ1 (Simulation fidelity):} To what extent can instruction-tuned LLMs simulate the \emph{distribution} of human responses in behavioural experiments, and where do they systematically deviate?
    \item \textbf{RQ2 (Opinion alignment and steering):} Whose opinions do LMs reflect, and how far can prompting towards a demographic group change this?
    \item \textbf{RQ3 (Generation):} How readily do current LLMs produce disinformation news articles, and how well do existing detectors recognise such articles?
    \item \textbf{RQ4 (Detection):} What do LLM-based approaches add to propaganda-technique detection compared with the systems of the SemEval-2020 shared task?
\end{enumerate}

\paragraph{Contribution.} Beyond summarising the four studies, this paper (i)~compares them along common dimensions (task, data, models, language, evaluation and main finding; Table~\ref{tab:comparison}); (ii)~places them in one conceptual framework (Figure~\ref{fig:framework}); (iii)~tests whether instruction tuning alone causes the hyper-accuracy distortion, using an open base/instruct model pair (Section~\ref{sec:experiment}); (iv)~derives two cross-study tensions, the \emph{generation--detection asymmetry} and the \emph{controllability gap}, that are not stated in any single study (Section~\ref{sec:comparative_synthesis}); and (v)~identifies open research gaps (Section~\ref{sec:governance}).
